\section{Conclusion and Limitations}
\label{sec:conclusion}

A commercial probe feed and a dashboard camera, combined in a global-local advisory, let ordinary vehicles prevent jams that the feed alone reports too late: the feed coordinates every vehicle on one network-wide plan without messages, and the camera lowers one vehicle's speed minutes before the feed's speed would warn. The system runs on US\$500 of commodity hardware, and in simulation it prevented the jam in 50 of 75 triggered runs and cut jam time by 97\% while serving 2.5\% more traffic than no control.

\noindent\textbf{Limitations.} No driver acted on an advisory, so the traffic results are simulated on one network, and they need most vehicles to follow: with half of them following, the system left more jam than no control after the largest burst. With no trigger the local layer reduced throughput by 6.9\%, and the feed gate that cuts this loss to 3.4\% leaves more jam after a stopped vehicle. The camera's distance scale is assumed, and the feed's delivery process was measured on one day.
